\documentclass[11pt, a4paper]{article}
\usepackage[utf8]{inputenc}
\usepackage[margin=1in]{geometry}
\usepackage{xcolor}
\usepackage{enumitem}
\usepackage{tcolorbox}
\usepackage{titlesec}

% 自定义颜色
\definecolor{njustblue}{RGB}{0, 51, 102}
\definecolor{highlightred}{RGB}{190, 0, 0}

% 标题设置
\title{\textbf{Academic Presentation Skills: Structure and Signposting}}
\author{Lecture Notes | Nanjing University of Science \& Technology}
\date{January 2026}

\begin{document}

\maketitle

\section{A Standard Academic Structure}

A successful academic presentation typically follows a logical 7-step progression:

\begin{enumerate}
    \item \textbf{Title \& Context:} Paper title, author’s name, affiliation. 
    \begin{itemize}
        \item \textit{Contexts:} Conference, course, seminar, workshop, lecture, etc.
    \end{itemize}
    
    \item \textbf{Introduction (Why this matters):} 
    \begin{itemize}
        \item Background (macro) and specific context.
        \item Research gap, problem, or question.
        \item Purpose, scope, focus, and significance of the study.
    \end{itemize}
    
    \item \textbf{Literature / Theoretical Framework (brief):}
    \begin{itemize}
        \item Key studies or theories; research gap; your position.
        \item Address: \textbf{what} has been studied, \textbf{how} it was approached, \textbf{what is missing}, and where \textbf{your study} stands.
    \end{itemize}
    
    \item \textbf{Methodology (How you did it):} 
    \begin{itemize}
        \item Data/materials, methods/approach, and justification (why this method?).
    \end{itemize}
    
    \item \textbf{Results / Analysis (\textcolor{highlightred}{Core part}):}
    \begin{itemize}
        \item Key findings only; use tables, graphs, and examples.
        \item \textbf{Highlight patterns}, do \textit{not} just present raw data.
    \end{itemize}
    
    \item \textbf{Discussion (What it means):}
    \begin{itemize}
        \item Interpretation of results; relation to research question; comparison with previous studies.
    \end{itemize}
    
    \item \textbf{Conclusion:}
    \begin{itemize}
        \item Main takeaways (2--3 points), contributions, limitations, and brief future research directions.
    \end{itemize}
\end{enumerate}

---

\section{Identifying Signposts}

\textbf{Definition:} Words or expressions indicating the connections between sentences and paragraphs.

\subsection{Functions and Applications}
Signposts enhance comprehension and help the audience make predictions. They are widely applied to:
\begin{itemize}
    \item Make a list / Show sequence
    \item Shift topics / Add information
    \item Provide examples / Clarify viewpoints
    \item Show cause \& effect
\end{itemize}

\subsection{Key Signposting Phrases}

\begin{tcolorbox}[colback=blue!5!white,colframe=blue!75!black,title=To Contrast \& Compare]
\textbf{Contrast:} but; yet; although; however; nonetheless; on the other hand; in contrast to; nevertheless. \\
\textbf{Compare:} similarly; likewise; in comparison with; similar to; correspondingly; equally. \\
\textit{Example:} \textbf{Similarly}, if you have food allergies, you should always ask what is in the dish.
\end{tcolorbox}

\begin{tcolorbox}[colback=red!5!white,colframe=red!75!black,title=To Give an Example]
for example; for instance; to illustrate; as an illustration; one/an example; as can be seen by; in other words; namely; specifically; a typical case in point is...; to be specific... \\
\textit{Example:} \textbf{A case in point} was the controversy provoked by the agitation for full payment.
\end{tcolorbox}

---

\section{Slide Design: Academic but Readable}

\begin{table}[h]
\centering
\begin{tabular}{|p{0.45\textwidth}|p{0.45\textwidth}|}
\hline
\textbf{\textcolor{blue}{Dos}} & \textbf{\textcolor{highlightred}{Don'ts}} \\ \hline
$\bullet$ One idea per slide & $\bullet$ Reading slides word for word \\
$\bullet$ Bullet points, not paragraphs & $\bullet$ Dense text \\
$\bullet$ Use figures instead of text where possible & $\bullet$ Decorative animations \\
$\bullet$ Consistent font and layout & \\ \hline
\end{tabular}
\end{table}

\end{document}